\documentclass[11pt]{article}
\usepackage{amsmath,amssymb,amsthm}
\usepackage[margin=1.1in]{geometry}
\usepackage{microtype}
\usepackage[hidelinks]{hyperref}
\setlength{\emergencystretch}{3em}

\newtheoremstyle{erratumstyle}{}{}{}{}{\bfseries}{.}{ }{#1 #2 (#3)}
\theoremstyle{erratumstyle}
\newtheorem{erratum}{Erratum}

\newcommand{\lgn}{\lg n}

\title{Errata and Counterexamples for\\
  \emph{Approximate Majority Analyses using\\
  Tri-molecular Chemical Reaction Networks}\\[2mm]
  {\large Condon, Hajiaghayi, Kirkpatrick, Ma\v{n}uch}\\[4mm]
  {\large found while formalizing the paper in Lean 4 / Mathlib}}
\author{Xiang Huang}
\date{}

\begin{document}
\maketitle

\begin{abstract}
These errata were found while producing a machine-checked formalization of the
paper in Lean 4 with Mathlib. Each entry identifies the exact printed passage,
gives either a short quotation or an exact mathematical transcription, and
states the correction; where a corrected theorem has been machine-checked we
name the corresponding Lean declaration. Every location below was checked
against the published Springer PDF of \emph{Natural Computing} 19 (2020),
249--270, DOI \href{https://doi.org/10.1007/s11047-019-09756-4}
{\texttt{10.1007/s11047-019-09756-4}}. We give both the printed publication
page and the one-based PDF page, counted from the article's title page.

The paper's entry and consensus conclusions are not refuted by this audit. One
additional persistence clause --- Theorem~4's preservation of relaxed consensus
for the following $n^{\gamma}$ interactions --- is false as printed. The
remaining items are proof gaps, omitted hypotheses or qualifiers, and
typographical errors. We do \emph{not} claim that every downstream use survives:
the false preservation clause has no corrected downstream proof at its printed
strength. The full-range tail of Lemma~12 / Theorem~5 is true, but the printed
proof does not establish it; a replacement state-dependent proof has been
machine-checked.

Classification scheme. (E) a printed statement is false; (P) a genuine proof
erratum with the statement still true; (U) an assumption, scope condition,
probability qualifier, or stronger conclusion is under-stated; (T) a typo or
typesetting omission with unambiguous intended meaning; and (O) \emph{open
after audit} --- the printed proof does not establish the stated parameter
range, and no counterexample or replacement proof is currently known.

The (O) label exists because E/P/U/T do not honestly cover a claim whose proof
fails while its statement is neither proved nor refuted: labelling such a claim
(P) would, by the definition above, assert that the statement is known true.
Exactly one defect is (E) --- Theorem~4's preservation clause, recorded first
below and refuted by an explicit fixed-adversary family. No defect remains in
class (O): the former open case, the full-range Lemma~12 / Theorem~5 tail, is
now classified (P). The replacement proof retains each competitor's actual
initial gap before summing the failure probabilities and yields the canonical
Lean declarations \texttt{Tri.Multi.lemma12\_unconditional} and
\texttt{Tri.Multi.theorem5}.
\end{abstract}

\section*{Notation}

Throughout, $n$ is the population size, $x$ and $y$ are the counts of the two
species with $x+y=n$, $\gamma \ge 1$ is the paper's confidence parameter, and
$\lg$ is the base-$2$ logarithm. A reaction event is \emph{productive} if it
changes the counts.

%--------------------------------------------------------------------
\begin{erratum}[Section 4.4 --- \texttt{creation-noise tolerance (P)}]
\textbf{Location.} \emph{Natural Computing} 19 (2020), p.~259
(Springer PDF p.~11), Sect.~4.4, first paragraph: the choice of the tolerance
used to control the imbalance between the two half-rate creations.

\medskip
\textbf{As printed.} The argument fixes the tolerance
\[
  d \;=\; \frac{\gamma n\lgn}{8\Delta}
\]
and asserts creation-noise control at $\exp(-\Theta(\gamma\lgn))$ uniformly on
the stated range $\sqrt{\gamma n\lgn}/2 < \Delta < n$.

\medskip
\textbf{Issue.} With a per-stage creation budget of order $n$, the resulting
Chernoff exponent is of order $d^2/n$, i.e.\ of order
$\gamma^2 n\lg^2 n/\Delta^2$. This is not $\Omega(\gamma\lgn)$ uniformly on the
printed range; it degrades quadratically as $\Delta$ grows and is already far
too small at the LOW end of the range.

\medskip
\textbf{Why.} At $n = 2^{20}$, $\gamma = 1$ (so $\gamma\lgn = 20$), with a
budget $15n$:
\[
  \begin{array}{c|c|c}
    \Delta & d & d^2/(2\cdot 15 n)\\\hline
    \sqrt{\gamma n\lgn}/2 & 1144.9 & 0.042\\
    n/16 & 40.0 & 1.1\times10^{-4}\\
    n/4 & 10.0 & \approx 0\\
    n/2 & 5.0 & \approx 0
  \end{array}
\]
Every entry is far below the required $\gamma\lgn = 20$. The defect is therefore
not confined to a large-$\Delta$ corner.

\medskip
\textbf{Fix.} Use stage-specific creation budgets and deviations rather than a
single global tolerance tied to $\Delta$. The formalization does exactly that:
\texttt{Tri.theorem2\_singleB} carries a per-stage $(D,H)$ pair, and the
underlying concentration lemma is in fact \emph{stronger} in temporal scope than
the printed fixed-count step --- \texttt{Tri/SingleBCreationTailY.lean} proves
the maximal stopped bound
\[
  \Pr\bigl[\exists t:\ C_t \le H \text{ and } CY_t - CX_t \ge D\bigr]
    \le \exp\!\left(-\frac{D^2}{2H}\right)
\]
with an explicit constant, where the paper states only an asymptotic
$\exp(-\Theta(\gamma\lgn))$. The theorem-level conclusion of the Single-B clause
is unaffected.
\end{erratum}

\begin{erratum}[Theorem 4 --- \texttt{preservation clause (E)}]
\textbf{Location.} \emph{Natural Computing} 19 (2020), pp.~261--262
(Springer PDF pp.~13--14), Theorem~4's second conclusion and its preservation
argument for the subsequent $n^{\gamma}$ interaction events.

\medskip
\textbf{As printed.} Under $\Delta_0 := 2x-n \ge \sqrt{\gamma n\lgn}$ and
$z < \Delta_0/16$, with probability $1-\exp(-\Omega(\gamma\lgn))$ a relaxed
$X$-consensus is reached within $O(\gamma n\lgn)$ interactions \emph{and is
preserved for the subsequent $n^{\gamma}$ interaction events}.

\medskip
\textbf{Why it is false.} The dynamics do not depend on the confidence
parameter $\gamma$, whereas the window length $n^{\gamma}$ does. With the
Byzantine count $z$ a fixed constant, escaping the buffer costs only
$n^{-O(z)}$ per attempt, and a window of unbounded polynomial degree
overwhelms it. Concretely, take
\[
  z = 2,\qquad \gamma = 17,\qquad (x,y,z) = (n-16,\,14,\,2),
\]
which lies exactly on the printed relaxed-consensus boundary $x = n-8z$ and
satisfies both hypotheses for all large $n$ (here $\Delta_0 = n-32$, and
$(n-32)^2 \ge 17\,n\lgn$ and $2 < (n-32)/16$).

Use the paper's own declared worst response, in which Byzantine molecules
always favour $Y$; this is a \emph{fixed, non-adaptive} adversary, so allowing
adaptive adversaries cannot rescue the claim. Partition the raw timeline into
blocks of $29n^2$ interactions. From any state with $0\le y\le 14$, one block
drives the process to $y=15$ --- i.e.\ out of the relaxed guard --- with
probability at least $c\,n^{-14}$, where $c = 1/(2\cdot 4^{14}\cdot 6^{14})$:
force $y$ down to $0$; nucleate $y=1$ from a triple $\{X,B,B\}$ whose two
Byzantine molecules both present $Y$, so that the reaction reads
$X+Y+Y\to 3Y$ and an honest $X$ becomes an honest $Y$; then take fourteen
prescribed up-steps, each costing at most a factor $1/(6n)$. There are
$\Theta(n^{15})$ such blocks in $n^{17}$ interactions, so
\[
  \Pr[\text{no exit by } n^{17}]
  \;\le\; \bigl(1-c\,n^{-14}\bigr)^{\Theta(n^{15})}
  \;\le\; \exp(-\Theta(n)) \;\longrightarrow\; 0,
\]
whereas the theorem claims an exit probability of
$\exp(-\Omega(\gamma\lgn))\to 0$.

\medskip
\textbf{Scale, stated honestly.} The constant is small: the crossover, where
$cn/29$ first exceeds $1$, is near $n \approx 1.2\times10^{21}$. Since the
theorem is asymptotic this is a legitimate refutation, and the crossover lies
far below the $n_0$ at which this development's other headline theorems are
even asserted.

\medskip
\textbf{Consequence.} Theorem~4's \emph{entry} conclusion is unaffected. The
preservation conclusion requires either a window tied to the Byzantine count
(rather than an unbounded polynomial degree), or a buffered guard with slack
growing in $\gamma$. The Lean development records the printed clause in
\texttt{Tri/Theorem4Statement.lean} as known-false and does not attempt it;
the corrected reached-by-deadline entry clause is machine-checked as
\texttt{Tri.Byzantine.theorem4\_entry :
Theorem4\_entry\_statement}, with failure at most $n^{-\gamma/136}$.

\medskip
\textbf{Related printed slips in the same passage.} The definition of relaxed
consensus by the \emph{equality} $x = n-8z$ cannot be preserved under any
dynamics and must be read as $x \ge n-8z$; the proof's ``$x \ge n-2z$ is
equivalent to $y < z$'' should read $y \le z$; and ``two doublings'' from
$M = 2z$ reaches the boundary $M = 8z$, so strict exit needs one further
upward step.
\end{erratum}

\begin{erratum}[Lemma 4 --- \texttt{proof}]
\textbf{Location.} \emph{Natural Computing} 19 (2020), p.~256
(Springer PDF p.~8), Lemma~4 proof, first displayed lower bound on the
productive-reaction success probability.

\medskip
\textbf{As printed.} After setting $d=\gamma\lg n$ and restricting to the
process before the first time $y>n/k+d$, the paper prints
\[
  \frac{x}{n}\;\ge\;1-\frac{n-kd}{kn}\;>\;1-\frac1k ,
\]
and concludes that each productive reaction succeeds with probability
$p>1-1/k$.

\medskip
\textbf{Issue.} Both the displayed inequality and its conclusion fail. From
$y\le n/k+d$ one gets $x=n-y\ge n-n/k-d$, hence
\[
  \frac xn\;\ge\;1-\frac1k-\frac dn\;=\;1-\frac{n+kd}{kn},
\]
so the numerator is $n+kd$, not $n-kd$: the printed line overstates $x/n$ by
$2d/n$. With the correct sign the bound reads $1-1/k-d/n<1-1/k$, so the
conclusion $p>1-1/k$ does not follow, and it is genuinely false at states
inside the band the proof claims to cover. The exact per-productive-event
success probability
\[
  p_y=\frac{\binom x2 y}{\binom x2 y+x\binom y2}=\frac{x-1}{n-2}
\]
does not repair it either. Numerical witnesses at the band's upper edge
$y=n/k+d$:
\[
  \begin{array}{c|c|c|c|c}
    k & n & d & p_y & 1-1/k\\\hline
    4 & 10^3 & 50 & 0.7004 & 0.7500\\
    8 & 10^3 & 50 & 0.8257 & 0.8750\\
    6 & 4096 & 60 & 0.8190 & 0.8333\\
    4 & 10^6 & 20 & 0.74998 & 0.75000
  \end{array}
\]
In each case $p_y<1-1/k$, so the claimed uniform floor does not hold on the
band.

\medskip
\textbf{Consequence.} Lemma 4's conclusion is nevertheless true. The
formalization (\texttt{Tri/PaperLemma4.lean}) proves it by a corrected
stopped argument: with $P=n/k$ the initial minority it uses the smaller safety
buffer $B=\lfloor P/8\rfloor$, on whose band the exact productive kernel does
have a uniform adaptive success floor; a Feller bound controls escape from the
band and an adapted Chernoff bound controls failure to reach consensus by the
paper's productive deadline. So this is a defect of the printed proof, not of
the statement.
\end{erratum}

%--------------------------------------------------------------------
\begin{erratum}[Corollary 2 --- \texttt{statement}]
\textbf{Location.} \emph{Natural Computing} 19 (2020), p.~256
(Springer PDF p.~8), displayed statement of Corollary~2.

\medskip
\textbf{As printed.}
``Since stage $s$ of phase 2 starts with $y \le n/2^s \le n/6$ and ends with
$y = \min\{n/2^{s+1},\ \gamma \lgn\}$, it completes properly, within at most
$3n/2^s$ productive reaction events, with probability at least
$1 - \exp(-\Theta(\gamma \lgn))$.''

\medskip
\textbf{Issue.} The $\min$ should be a $\max$. As printed, the stated
end-value is unattainable in either regime.

\medskip
\textbf{Why.} The paper's own definition of a phase-2 stage (L825) is
unambiguous:

\begin{quote}
``A typical stage $s$ in phase 2 starts with $y \le n/2^s$ and continues while
$n/2^{s-1} \ge y > n/2^{s+1}$. It completes properly if $y \le n/2^{s+1}$.''
\end{quote}

So the stage's own target is $y \le n/2^{s+1}$, while phase 2 as a whole
terminates at its floor $y \le \gamma\lgn$ (L811). Since $y$ is
non-increasing across the relevant events, a stage that may terminate for
either reason ends at whichever bound is reached \emph{first}, i.e.\ at the
\emph{larger} of the two. Concretely:

\begin{itemize}
\item if $\gamma\lgn > n/2^{s+1}$, phase 2 has already terminated before $y$
  descends to $n/2^{s+1}$, so the printed $\min = n/2^{s+1}$ is never reached
  within the phase;
\item if $\gamma\lgn < n/2^{s+1}$, the stage completes at the dyadic threshold
  $n/2^{s+1}$, not at $\gamma\lgn$, so the printed $\min = \gamma\lgn$ is
  again not the end-value.
\end{itemize}

Either way the printed $\min$ is wrong, and $\max$ is correct.

\medskip
\textbf{Fix.} Replace $\min$ by $\max$. Nothing downstream changes: the proof
of Corollary~2 only uses that $y$ stays above the live-band floor during the
stage, which is exactly what the $\max$ form asserts.

The Lean development sidesteps the issue entirely by formalizing the
\emph{stage definition} (L825), which involves no $\min$ or $\max$: writing
$P = n/2^s$ and $B = n/2^{s+1}$, the stage runs from $y \le P$, escapes if
$y > 2P$, and completes at $y \le B$, within $3P$ productive events.
\end{erratum}

%--------------------------------------------------------------------
\begin{erratum}[Lemma 6 --- \texttt{threshold}]
\textbf{Location.} \emph{Natural Computing} 19 (2020), p.~260
(Springer PDF p.~12), displayed statement of Lemma~6.

\medskip
\textbf{As printed.} Under the unequal-rate hypothesis, the gap
``increases to $\min\{2\Delta,n\}$'' within $5n$ productive reaction events.

\medskip
\textbf{Issue.} Classification: (T). Productive reactions change the gap by
$2$, and all reachable gaps have the same parity as $n$. Thus the displayed
target cannot be read as exact equality in every population size. A concrete
witness is $n=9$, initial split $(x,y)=(5,4)$, and $\Delta=1$: the reachable
gaps are $\{1,3,5,7,9\}$, so the literal target $2$ is unreachable.

\medskip
\textbf{Fix.} Read the conclusion as reaching the first reachable gap at least
the threshold, i.e.\ $G_t\ge \min\{2\Delta,n\}$. This is the intended
threshold statement and is what the proof supports.
\end{erratum}

%--------------------------------------------------------------------
\begin{erratum}[Lemma 12 --- \texttt{gap (P)}]
\textbf{Location.} \emph{Natural Computing} 19 (2020), p.~263
(Springer PDF p.~15), Lemma~12 proof, the type-(a) estimate immediately after
the four stopping alternatives.

\medskip
\textbf{As printed (statement).}
``Consider a point in the computation where $x - z_{\max} = \Delta \ge
\sqrt{\gamma n \lgn}$. Let $m \le n/(\gamma\lgn)$. Then
$x - z_{\max} \ge \min\{2\Delta,\ n\}$ within $\Theta(n)$ productive reaction
events, with probability $1 - \exp(-\Omega(\gamma\lgn))$.''

\medskip
\textbf{As printed (the step at issue, in the proof).}
``By Lemma 11, completion of type (a) occurs with probability
$m\exp(-\Omega(\gamma\lgn)) = \exp(-\Omega(\gamma\lgn))$ since $m < n$.''

\medskip
\textbf{Issue.} Classification: (P). The displayed absorption of the factor $m$ into the
$\Omega(\cdot)$ is not justified by $m < n$ alone. It additionally requires the
constant hidden in the exponent to exceed $\ln 2$, which the paper never
establishes and which the natural constants coming out of the argument do not
satisfy.

\medskip
\textbf{Why.} Write the exponent explicitly: Lemma 11 supplies a failure mass
of the form $m\exp(-c\,\gamma\lgn)$ for some constant $c>0$ determined by the
argument. Using only $m \le n = 2^{\lgn}$,
\[
  m\exp(-c\,\gamma\lgn)\ \le\ \exp\bigl((\ln 2)\lgn - c\,\gamma\lgn\bigr)
  \ =\ \exp\bigl(-(c\gamma - \ln 2)\lgn\bigr).
\]
For this to be $\exp(-\Omega(\gamma\lgn))$ one needs $c\gamma - \ln 2 \ge
c'\gamma$ for some $c'>0$, i.e.
\[
  \gamma\,(c-c')\ \ge\ \ln 2 .
\]
Since the lemma is asserted ``for any constant $\gamma \ge 1$'', this forces
$c > \ln 2 \approx 0.693$. The constant that the machine-checked version of
this argument actually produces is $c = 1/82944$, so for that constant the
absorption fails outright for every $\gamma$ in a very large range. The
paper's own hypothesis $m \le n/(\gamma\lgn)$ helps only by a
$\log$ factor and does not change the conclusion.

\medskip
\textbf{Fix.} The earlier formal wrapper exposed the exact additional
condition required by the printed coarse union-bound route:
\[
  288m + 144\ \le\ \exp\!\left(\frac{\gamma\lgn}{165888}\right),
\]
but the replacement proof does not use that route. For a competitor with
initial gap $g_Y$, it retains both the Gaussian decay at the minimum gap
$\Delta$ and the additional decay proportional to $g_Y-\Delta$. Competitors
with population comparable to $x_0$ are then counted by total population
rather than by $m$, while smaller competitors are absorbed by the additional
decay. The resulting 144-substage error is at most
\[
  432\exp\!\left(-\frac{\gamma\lgn}{663552}\right)
\]
throughout the printed range $m\gamma\lgn\le n$ and
$\gamma n\lgn\le\Delta^2$. This is
\path{Tri.Multi.lemma12_unconditional}.

\medskip
\textbf{Remark.} The printed statement is true. The error is in discarding the
competitor-dependent decay before summing and then trying to absorb the
resulting artificial factor $m$.
\end{erratum}

%--------------------------------------------------------------------
\begin{erratum}[Lemma 12 --- \texttt{proof}]
\textbf{Location.} \emph{Natural Computing} 19 (2020), p.~264
(Springer PDF p.~16), Lemma~12 proof, the paragraph beginning ``Next, consider
completions of type (c),'' where progress against a fixed species is followed
by an invocation of Lemma~11.

\medskip
\textbf{As printed.} The proof first obtains, for a fixed species $Y$, a new
large $X$--$Y$ gap, and then feeds that conclusion into Lemma~11.

\medskip
\textbf{Issue.} Classification: (P). Lemma~11 assumes a \emph{global}
plurality gap $x-z_{\max}\ge \Delta'$. A large gap against one fixed $Y$ at
that species' individual hitting time does not by itself establish the global
premise against every competitor. This is logically separate from the
species-factor absorption error above.

\medskip
\textbf{Fix.} The replacement proof uses one synchronized global stopping
process. Its safety and progress potentials are summed over all competitors
while retaining each actual initial pair gap. Consequently every terminal
failure is charged once, and no fixed-$Y$ conclusion is passed to the global
Lemma~11. This proof is implemented by
\texttt{productiveProperStage\_progress\_state},
\texttt{productiveProperStageStateError\_le\_headline}, and the public
\texttt{lemma12\_unconditional}.
\end{erratum}

%--------------------------------------------------------------------
\begin{erratum}[Lemma 12 --- \texttt{terminal line}]
\textbf{Location.} \emph{Natural Computing} 19 (2020), p.~264
(Springer PDF p.~16), final sentence of the Lemma~12 proof.

\medskip
\textbf{As printed.} The proof ends with the uncapped conclusion
$x-z_{\max}\ge 2\Delta$.

\medskip
\textbf{Issue.} Classification: (T). The uncapped inequality is impossible
when $\Delta>n/2$. This is a terminal typesetting omission, since the lemma's
own statement already uses the capped target.

\medskip
\textbf{Fix.} The final line should read
$x-z_{\max}\ge \min\{2\Delta,n\}$, with the terminal branch interpreted as
consensus.
\end{erratum}

%--------------------------------------------------------------------
\begin{erratum}[Lemma 13 --- \texttt{exponent}]
\textbf{Location.} \emph{Natural Computing} 19 (2020), p.~264
(Springer PDF p.~16), final sentence of the Lemma~13 proof.

\medskip
\textbf{As printed.} The proof says that $\lambda mn$ raw interactions, with
a productive-reaction floor of order $\Theta(1/m)$, give failure
$\exp(-\Theta(n/m))$.

\medskip
\textbf{Issue.} Classification: (U). The Chernoff mean for the productive
count is actually $\Theta(n)$, so the constant-fraction lower-tail calculation
gives $\exp(-\Theta(n))$. The printed exponent understates the calculation.

\medskip
\textbf{Fix.} State the stronger $\exp(-\Theta(n))$ clock tail. The weaker
reading $\exp(-\Omega(n/m))$ is harmless for the downstream use.
\end{erratum}

%--------------------------------------------------------------------
\begin{erratum}[Theorem 3 --- \texttt{proof candidate}]
\textbf{Location.} \emph{Natural Computing} 19 (2020), p.~261
(Springer PDF p.~13), proof of Theorem~3, parameter-algebra paragraph
immediately following the theorem statement.

\medskip
\textbf{As printed.} In the proof's parameter algebra, a displayed
intermediate line asserts $c<\alpha$ from assumptions including
$0<\xi<\alpha<1$.

\medskip
\textbf{Issue.} Classification: audit-flagged (P) candidate, not
independently re-verified here. The implication $c<\alpha$ is false in
general; for example, the audit reports that $\alpha=0.1$ and $\xi=0.05$ give
$c\approx0.859$. The line appears unnecessary for the proof, whose continuation
needs $c<1$ and the later $k,\beta$ inequalities.

\medskip
\textbf{Fix.} Delete the false intermediate inequality, or replace it by the
actually used inequalities after verifying the local derivation. This entry is
recorded as a candidate exactly as flagged by the audit.
\end{erratum}

%--------------------------------------------------------------------
\begin{erratum}[Lemma 10 --- \texttt{finite size}]
\textbf{Location.} \emph{Natural Computing} 19 (2020), p.~262
(Springer PDF p.~14), Lemma~10 proof.

\medskip
\textbf{As printed.} From $k\ge4$ and $x<n-2z$, the proof derives the displayed
effective-rate inequality under the paper's informal ``for $n$ sufficiently
large'' convention.

\medskip
\textbf{Issue.} Classification: (P) side condition. The finite statement is
only asymptotically valid. A uniform sufficient condition is
\[
  \gamma\lg n \le \frac{3n}{68}.
\]
The common standing assumption $6\gamma\lg n\le n$, equivalently
$\gamma\lg n\le 11.33n/68$, is weaker and therefore does not imply this
Lemma~10 instance.

\medskip
\textbf{Fix.} Carry the strict integer form
$68\gamma\lg n+1\le 3n$, or an equivalent explicit hypothesis, in every
formal use of Lemma~10. The corrected Theorem~4 entry declaration
\texttt{Tri.Byzantine.theorem4\_entry} carries exactly this side condition.
\end{erratum}

%--------------------------------------------------------------------
\begin{erratum}[Theorem 4 --- \texttt{strictness}]
\textbf{Location.} \emph{Natural Computing} 19 (2020), p.~262
(Springer PDF p.~14), Theorem~4 proof immediately downstream of Lemma~10.

\medskip
\textbf{As printed.} The proof treats $x\ge n-2z$ as implying $y<z$.

\medskip
\textbf{Issue.} Classification: (T). Since $x+y+z=n$, the inequality
$x\ge n-2z$ is equivalent to $y\le z$, not to $y<z$.

\medskip
\textbf{Fix.} Replace the strict inequality by $y\le z$ and adjust the
stopping predicate accordingly.
\end{erratum}

%--------------------------------------------------------------------
\begin{erratum}[Lemmas 17--19 --- \texttt{probability qualifiers}]
\textbf{Location.} \emph{Natural Computing} 19 (2020), pp.~266--267
(Springer PDF pp.~18--19): the displayed statements of Lemmas~17 and 18 on
p.~266, and Lemma~19 on p.~267.

\medskip
\textbf{As printed.} The displayed lemmas have no probability qualifier.

\medskip
\textbf{Issue.} Classification: (T), but formalization-critical. Read
literally, the statements are deterministic, which is false: adverse finite
reaction or activation sequences have positive probability. The proofs and
downstream uses are high-probability statements.

\medskip
\textbf{Fix.} Attach the intended high-probability qualifier to each lemma,
conditional on its stated start-of-stage hypotheses and inherited epidemic
events:
\begin{itemize}
\item Lemma~17: from the stated range
$n/(\gamma n\lg n)^{1/5}\le a<n/d^*$ and
$E_Y(a)\le 14c^*\sqrt{\gamma a\lg n}$, the conclusion
$E_Y(2a)\le 14c^*\sqrt{\gamma(2a)\lg n}$ holds with probability
$1-\exp(-\Omega(\gamma\lg n))$.
\item Lemma~18: under the preceding Lemma~17 start-of-stage bound, the stage
$a=n/d^*\to 2n/d^*$ creates active $X$-excess at least
$2Q$, where $Q=\sqrt{\gamma n\lg n}$, with probability
$1-\exp(-\Omega(\gamma\lg n))$.
\item Lemma~19: if the active $X$-excess is at least $2Q$ at
$2n/d^*$ activations, then at full activation the active $X$-excess is at
least $Q$ with probability $1-\exp(-\Omega(\gamma\lg n))$.
\end{itemize}
\end{erratum}

%--------------------------------------------------------------------
\begin{erratum}[Lemma 18 --- \texttt{stopping}]
\textbf{Location.} \emph{Natural Computing} 19 (2020), p.~266
(Springer PDF p.~18), Lemma~18 proof, in the sentence preceding estimates
(i)--(iii) and in the subsequent stopped rate estimate.

\medskip
\textbf{As printed.} With $Q=\sqrt{\gamma n\lg n}$, the proof asserts that if
the excess of active reactions of type (2) over type (1) is at most $16Q$, then
the stage completes without the active $Y$-excess ever exceeding $30Q$.

\medskip
\textbf{Issue.} Classification: (P), correcting our own prior classification
of this point as a non-erratum. Let $E_t$ be active $Y$-excess and $R_t$ be
the running reaction excess. On the preceding high-probability events,
$E_0\le 14Q$ and the activation-prefix contribution is at most $Q$, so
\[
  E_t \le 15Q + 2R_t.
\]
Crossing the $30Q$ barrier therefore forces only
$R_t>15Q/2=7.5Q$, not $R_t>16Q$. For example, $R_t=10Q\le16Q$ permits
$E_t=35Q>30Q$. The stopping-time phrasing in step~(ii) correctly makes the
local rate estimate non-circular, but it does not rule out exiting through the
$30Q$ barrier before the claimed endpoint reaction-excess event.

\medskip
\textbf{Fix.} Add a maximal stopped-walk bound
\[
  \max_{t\le\tau} R_t \le \frac{15Q}{2}
\]
with failure $\exp(-\Omega(\gamma\lg n))$, where $\tau$ is the minimum of the
stage end and the first $30Q$ crossing. The paper's own rate estimate and
active-reaction-count bound supply this estimate. It is stronger than the
endpoint bound needed later, since $15Q/2<16Q$.
\end{erratum}

%--------------------------------------------------------------------
%--------------------------------------------------------------------
\begin{erratum}[Heavy-B --- \texttt{exact emulation claims (P)}]
\textbf{Location.} \emph{Natural Computing} 19 (2020), p.~258
(Springer PDF p.~10), Sect.~4.1, ``Correctness of the Heavy-B emulation,''
last paragraph.

\medskip
\textbf{As printed.} With $\hat x=x+b$ and $\hat y=y+b$, the paper says that
``the relative probability of reaction (1') to that of (2'), equals the
relative probability of reaction (1) to that of (2) in the tri-molecular
CRN,'' and concludes that the evolution of $\hat x,\hat y$ ``exactly matches''
that of $x,y$ in Tri.

\medskip
\textbf{Issue.} The informal verb ``emulate'' is not itself the error. The
problem is the literal pair of claims used in the Section~4.1 correctness
argument: the reaction odds are asserted to be equal, and the virtual process
is asserted to match Tri exactly. Neither is true in general. This is a defect
in the printed proof, not a refutation of the Heavy-B clause of Theorem~2.

\medskip
\textbf{Why.} In a Heavy-B state $(x,y,b)$, conditioned on a resolution
reaction, the favorable and adverse propensities have ratio
\[
  xb:yb=x:y.
\]
At the virtual Tri state $(\hat x,\hat y)=(x+b,y+b)$, the two productive
tri-molecular propensities have ratio
\[
  \binom{\hat x}{2}\hat y:
    \hat x\binom{\hat y}{2}
  = (\hat x-1):(\hat y-1)
  = (x+b-1):(y+b-1).
\]
Equality would require
\[
  x(y+b-1)=y(x+b-1),
\]
or equivalently
\[
  (b-1)(x-y)=0.
\]
Thus the odds agree only when $b=1$ or $x=y$. For example,
$(x,y,b)=(3,1,3)$ gives Heavy-B odds $3:1$, whereas the projected Tri odds
are $5:3$. The raw interaction laws also use different entity populations:
Heavy-B samples from $x+y+b$ entities, while the virtual population has
$x+y+2b$ entities. These two obstructions are machine-checked in
\texttt{Tri/HeavyBEmulation.lean}: \texttt{Tri.heavy\_tri\_odds\_agree\_iff}
proves the exact equality criterion, and
\texttt{Tri.heavyToDouble\_entity\_count} proves the entity-count mismatch
whenever $b>0$.

\medskip
\textbf{Fix.} Replace the exact-match argument by the one-sided comparison
actually needed in the majority region: if $x\ge y$ and $b\ge1$, the Heavy-B
resolution skeleton is at least as biased toward $X$ as productive Tri, with
equality only in the two degenerate cases above. The proof must then combine
this stochastic bias comparison with a Heavy-B-specific resolution clock; it
cannot be obtained by black-box transport along an exact Markov lumping. The
Heavy-B theorem itself is proved directly as
\texttt{Tri.theorem2\_heavyB}.
\end{erratum}

%--------------------------------------------------------------------
\begin{erratum}[Single-B --- \texttt{creation-imbalance concentration (P)}]
\textbf{Location.} \emph{Natural Computing} 19 (2020), p.~259
(Springer PDF p.~11), Sect.~4.4, Single-B phase-1 analysis, paragraph beginning
``The probability that reaction $(0'_x)$ occurs is the same as the probability
that reaction $(0'_y)$ occurs.''

\medskip
\textbf{As printed.} The paper observes that the two creation reactions
$(0'_x)$ and $(0'_y)$ are equally likely and applies Lemma~2 to the number of
$(0'_x)$ reactions among $kn$ creation reactions. It then uses the resulting
endpoint count to conclude that the virtual gap $\hat x-\hat y$ does not
\emph{become} smaller than the stated lower barrier during the stage.

\medskip
\textbf{Issue.} Lemma~2 supplies concentration at a fixed endpoint, whereas the
needed assertion is maximal: the creation-orientation imbalance must remain
inside a barrier at every prefix before the stage terminates. An endpoint bound
does not exclude a trajectory that crosses the adverse barrier and later
returns. The printed argument also does not justify treating the
history-dependent sequence of creation times as an independent sequence of
trials.

\medskip
\textbf{Why.} Let $C_X$ and $C_Y$ count the two creation orientations. Their
masses are equal at every physical state, so $C_X-C_Y$ is an adapted fair walk
when indexed by creation events. Fairness alone does not convert a terminal
estimate into an ever-crossing estimate. The appropriate exponential process
on the full seven-event chain is
\[
  G_t=
  \exp\!\left(
    \lambda(C_X-C_Y)
    -\frac{\lambda^2}{2}(C_X+C_Y)
  \right).
\]
Inert and resolution reactions leave $G_t$ unchanged. The two creation
reactions multiply it by
$\exp(\lambda-\lambda^2/2)$ and
$\exp(-\lambda-\lambda^2/2)$ with equal mass, and
\[
  \frac{e^\lambda+e^{-\lambda}}2
  \le e^{\lambda^2/2}.
\]
Hence $G_t$ is a supermartingale without an independence assumption. If the
number of creations is at most $H$, Ville's inequality gives
\[
 \Pr\!\left[
   \exists t:\ C_X(t)-C_Y(t)\ge D,
   \ C_X(t)+C_Y(t)\le H
 \right]
 \le \exp\!\left(-\frac{D^2}{2H}\right),
\]
and the sign-reversed potential gives the opposite tail.

\medskip
\textbf{Fix.} Replace the endpoint invocation of Lemma~2 by the stopped
maximal estimate above, retain an explicit creation-count budget, and combine
both creation tails with the corrected resolution-gap process before
recovering the physical gap. The formalization proves the equal-mass input as
\texttt{Tri.singleCompPMF\_creation\_masses\_eq}, the full-kernel one-step
supermartingale as \texttt{Tri.creationG\_step}, the two maximal bounds as
\texttt{Tri.creation\_tail\_maximal} and
\texttt{Tri.creationY\_tail\_maximal}, and their optimized forms as
\texttt{Tri.creation\_tail} and \texttt{Tri.singleCreationY\_tail}. The
complete Single-B clause is \texttt{Tri.theorem2\_singleB}.
\end{erratum}

%--------------------------------------------------------------------
\begin{erratum}[Lemma 8 --- \texttt{adaptive-policy comparison (P)}]
\textbf{Location.} \emph{Natural Computing} 19 (2020), p.~261
(Springer PDF p.~13), Lemma~8 and its proof in the Byzantine section.

\medskip
\textbf{As printed.} For $0<\ell\le n-x-z$, Lemma~8 says that
\[
  \pi^*((x,y,z),(x+\ell,y-\ell,z))
\]
is minimized by the paper's fixed Byzantine response. The proof says it
suffices to observe
\[
 \pi^*((x,y,z),q)
 \le
 \pi^*((x+1,y-1,z),q),
 \qquad q=(x+\ell,y-\ell,z),
\]
because every path from the first state to $q$ passes through the second.

\medskip
\textbf{Issue.} The displayed inequality compares two starting states under
one fixed transition law. The lemma's conclusion compares different Byzantine
response policies, including responses that may depend on the complete past.
Starting one step closer to a target under a fixed policy does not establish
that one policy minimizes the target-hitting probability.

\medskip
\textbf{Why.} Skip-freeness proves only the stated spatial monotonicity: for a
fixed policy, every successful path from $x$ to $x+\ell$ first visits $x+1$.
A policy comparison needs a Bellman argument. For each finite horizon, the
continuation value---the probability of reaching the upper target in the
remaining steps---must first be proved monotone in the current $X$ count. The
pointwise worst-response inequality can then be applied to this monotone
continuation value, followed by backward induction over the horizon and over
the complete history. The eventual hitting comparison is obtained by taking
the increasing supremum over finite horizons. Those policy-comparison steps
are absent from the printed proof.

\medskip
\textbf{Fix.} Replace the proof by finite-horizon dynamic programming and then
pass to the unbounded hitting event. This repair is machine-checked as
\path{Tri.Byzantine.lemma8}. Its finite-horizon core
\path{lemma8_finite_horizon} proves Bellman domination against every
history-dependent strategy; nearest-neighbour support identifies the upper
Bellman target with the exact configuration, and a supremum over all finite
horizons gives the paper's unbounded $\pi^*$ comparison. The downstream
potential form remains available as
\path{controlledLaw_antitone_expect_le_pow} and
\path{controlledLaw_antitone_super}.
\end{erratum}

%--------------------------------------------------------------------
\begin{erratum}[Lemma 11 --- \texttt{ratio inequalities (P)}]
\textbf{Location.} \emph{Natural Computing} 19 (2020), p.~263
(Springer PDF p.~15), Lemma~11 proof, the two-party and third-party
backsliding calculations.

\medskip
\textbf{As printed.} In the two-party calculation the paper sets
$p=1/2+\Delta/(4n)$ and prints
\[
 \left(\frac{1-p}{p}\right)^{\Delta/8}
 < \left(1-\frac{\Delta}{n}\right)^{\Delta/8}.
\]
In the third-party calculation it sets
$p=1/2+\Delta/(8n)$ and prints
\[
 \left(\frac{1-p}{p}\right)^{\Delta/4}
 < \left(1-\frac{\Delta}{2n}\right)^{\Delta/4}.
\]

\medskip
\textbf{Issue.} Both displayed comparisons are false: before taking powers,
the left-hand ratio is strictly \emph{larger} than the printed linear factor.
These two displayed comparisons are the entirety of the ``false ratio
inequalities'' defect. No stronger defect in Lemma~11 is asserted: the
paper's two-mode decomposition is legitimate, and its stated
$m\exp(-\Omega(\Delta^2/n))$ conclusion survives after repairing these two
estimates.

\medskip
\textbf{Why.} In the two-party case,
\[
 \frac{1-p}{p}=\frac{2n-\Delta}{2n+\Delta},
\]
and
\[
 \frac{2n-\Delta}{2n+\Delta}
 -\left(1-\frac{\Delta}{n}\right)
 =\frac{\Delta^2}{n(2n+\Delta)}>0.
\]
For example, $n=100$, $\Delta=10$ gives
$190/210\approx0.90476>0.9$. In the third-party case,
\[
 \frac{1-p}{p}=\frac{4n-\Delta}{4n+\Delta},
\]
and
\[
 \frac{4n-\Delta}{4n+\Delta}
 -\left(1-\frac{\Delta}{2n}\right)
 =\frac{\Delta^2}{2n(4n+\Delta)}>0.
\]
For $n=100$, $\Delta=10$ this is
$390/410\approx0.95122>0.95$.

\medskip
\textbf{Fix.} Use the valid inequality
\[
  \frac{1-t}{1+t}\le e^{-2t}
  \qquad (0\le t<1).
\]
With $t=\Delta/(2n)$ in the two-party case,
\[
 \left(\frac{1-p}{p}\right)^{\Delta/8}
 \le \exp\!\left(-\frac{\Delta^2}{8n}\right).
\]
With $t=\Delta/(4n)$ in the third-party case,
\[
 \left(\frac{1-p}{p}\right)^{\Delta/4}
 \le \exp\!\left(-\frac{\Delta^2}{8n}\right).
\]
Union over the two modes and the $m-1$ competitors therefore preserves the
paper's asymptotic form $m\exp(-\Omega(\Delta^2/n))$. The Lean development
avoids the two erroneous linear comparisons and proves the all-time stopped
estimate directly by a geometric supermartingale. The core theorem is
\path{Tri.Multi.productivePairGapEver_half_failure_exp}; its two
paper-facing wrappers are
\path{Tri.Multi.lemma11_twoPartyBacksliding} and
\path{Tri.Multi.lemma11_thirdPartyBacksliding} in
\path{Tri/PaperLemma11.lean}, with the explicit envelope
\[
  m\exp\!\left(-\frac{\Delta^2}{18n}\right).
\]
\end{erratum}

\section*{Non-errata: differences that are ours, not the paper's}

For the record, the following divergences between the Lean development and the
paper are \emph{our} choices or \emph{our} earlier mistakes, and are
deliberately not listed above.

\begin{itemize}
\item Phase 0 of Theorem 1(a) is proved with a concave symmetric Lyapunov
  supermartingale rather than the paper's foray/checkpoint combinatorics. Both
  are correct; ours was easier to formalize.
\item Corollary 3's proof omits a monotonicity step: the productive-clock
  estimate may be proved at the boundary minority size $y=\gamma\lg n$ and
  then pushed to starts with smaller minority. This is a true omitted step, not
  an erratum, because the statement remains correct. The formalization records
  the missing comparison as
  \path{Tri.corollary3_boundary_success_le_entry_success}, built from
  the reusable Bellman comparison
  \path{Tri.productive_hitProb_mono_start_of_upper}.
\item An early version of our Corollary 3 wrapper stated a raw-interaction
  result under the productive-clock name. That was our labelling error, now
  corrected, and is not a defect of the paper.

\item Lemma 18's deterministic endpoint ledger closes with
  \emph{zero} slack. The paper's own arithmetic is
  $48\sqrt{\gamma n\lgn} - 14\sqrt{\gamma n\lgn} - 32\sqrt{\gamma n\lgn}
   = 2\sqrt{\gamma n\lgn}$,
  exactly the required gap. A formalization has no room to lose a constant
  anywhere in this lemma.
\end{itemize}

\end{document}
